\documentclass[12pt]{article}

\usepackage[margin=1in]{geometry}
\usepackage[T1]{fontenc}
\usepackage[utf8]{inputenc}
\usepackage{amsmath,amssymb}
\usepackage{booktabs}
\usepackage{setspace}
\usepackage[hidelinks]{hyperref}
\usepackage{graphicx}
\usepackage{siunitx}
\usepackage{caption}

\doublespacing
\graphicspath{{figures/}{./}}

\sisetup{
  group-separator = {,},
  group-minimum-digits = 4,
  table-number-alignment = center
}

\renewcommand{\thetable}{S\arabic{table}}
\renewcommand{\thefigure}{S\arabic{figure}}

\title{Supplementary Material\\Pooled angular profiling of contextual separation in implicit-association latency data}
\author{}
\date{}

\begin{document}

\maketitle

\section*{Supplementary Results}

\begin{table}[htbp]
\centering
\caption{Dataset-level summary of validation corpora.}
\label{tab:supp-dataset-summary}
\small
\setlength{\tabcolsep}{4pt}
\begin{tabular}{@{}p{3.1cm}p{3.2cm}p{2.5cm}p{2.5cm}p{3.0cm}@{}}
\toprule
Dataset & Source & Analytic sample & Critical blocks & Analytic role \\
\midrule
Gender--Science IAT & Project Implicit public archive & 141{,}329 participants & Blocks 3, 4, 6, and 7 & Primary public reference corpus \\
Sexuality IAT & Project Implicit public archive & 207{,}366 participants & Blocks 3, 4, 6, and 7 & Matched public social-category comparison \\
Age IAT & Project Implicit OSF archive & 156{,}005 participants after critical-block latency cleaning & Blocks 3, 4, 6, and 7 & Additional public social-category comparison \\
Life-satisfaction IAT & Independent Testable study & 180 participants & Blocks 3 and 5 & External self-referential comparison \\
\bottomrule
\end{tabular}

\begin{flushleft}
\footnotesize\emph{Note}. IAT = Implicit Association Test; OSF = Open Science Framework. The public IAT datasets preserve the standard group-based Project Implicit architecture, whereas the life-satisfaction IAT differs in platform, protocol, language, sample size, stimulus format, and construct target.
\end{flushleft}
\end{table}

\subsection*{Robustness to binning resolution}

To assess sensitivity to temporal discretization, the full profiling pipeline was repeated using 4, 6, and 8 quantile bins per participant. Across all three resolutions, the pooled profile remained concentrated in the low-angle regime. The posterior mode shifted from \(22.75^\circ\) at 4 bins to \(17.25^\circ\) at 6 bins and \(15.50^\circ\) at 8 bins, showing stable regime-level behavior with measurable discretization sensitivity.

\begin{table}[htbp]
\centering
\caption{Pooled profile-posterior summary by binning resolution.}
\label{tab:supp-posterior-bins}
\small
\begin{tabular}{@{}cccccccc@{}}
\toprule
Bins & Mean $\theta$ & SD & MAP & HDI low & HDI high & Posterior entropy & $\hat{\sigma}$ at MAP \\
\midrule
4 & 22.901 & 0.623 & 22.750 & 21.750 & 24.000 & 2.330 & 0.443 \\
6 & 17.292 & 0.483 & 17.250 & 16.500 & 18.250 & 2.077 & 0.630 \\
8 & 15.566 & 0.423 & 15.500 & 14.750 & 16.250 & 1.944 & 0.697 \\
\bottomrule
\end{tabular}
\end{table}

\subsection*{Bootstrap stability}

To assess the stability of the pooled estimate, we drew 100 bootstrap subsamples of $N=30{,}000$ participants from the full dataset and reran the profiling procedure for each subsample. The pooled MAP was $17^\circ$ in every replicate, indicating that the empirical low-angle solution is highly stable to resampling at large $N$.

\begin{figure}[htbp]
\centering
\includegraphics[width=0.72\textwidth]{ext_bootstrap_distribution.png}
\caption{Bootstrap stability of the pooled estimate. Each replicate resamples 30,000 participants and recomputes the pooled MAP.}
\label{fig:supp-bootstrap}
\end{figure}

\subsection*{Generic recovery diagnostics}

These simulations are included as diagnostics of the estimator's generic behavior across the full angular range. They are not the primary basis for interpreting the empirical low-angle regime, which is addressed separately below using matched-domain calibration under empirically realistic coefficient and noise structures.

\begin{table}[htbp]
\centering
\caption{Simulation-based recovery over the full profiling grid.}
\label{tab:supp-recovery}
\small
\begin{tabular}{@{}cccccccc@{}}
\toprule
True $\theta$ & Replicates & Mean MAP & Bias & RMSE & Median SD & 94\% coverage & Median HDI width \\
\midrule
5.000   & 100 & 41.743  & 36.743  & 40.930 & 14.394 & 0.000 & 50.125 \\
10.000  & 100 & 39.763  & 29.763  & 33.412 & 12.658 & 0.000 & 44.500 \\
15.000  & 100 & 42.363  & 27.363  & 31.090 & 12.451 & 0.000 & 44.500 \\
17.250  & 100 & 43.173  & 25.923  & 29.044 & 11.810 & 0.000 & 41.875 \\
20.000  & 100 & 41.613  & 21.613  & 24.666 & 10.752 & 0.010 & 38.625 \\
25.000  & 100 & 42.030  & 17.030  & 19.739 & 8.620  & 0.500 & 31.375 \\
30.000  & 100 & 39.915  & 9.915   & 11.987 & 6.995  & 0.800 & 25.625 \\
40.000  & 100 & 42.610  & 2.610   & 5.964  & 5.172  & 0.910 & 19.250 \\
60.000  & 100 & 59.948  & -0.053  & 2.524  & 2.356  & 0.910 & 8.750 \\
90.000  & 100 & 89.988  & -0.013  & 0.962  & 0.885  & 0.940 & 3.250 \\
120.000 & 100 & 119.990 & -0.010  & 0.494  & 0.438  & 0.950 & 1.500 \\
\bottomrule
\end{tabular}
\end{table}

\begin{figure}[htbp]
\centering
\includegraphics[width=0.95\textwidth]{ext_recovery_plots.png}
\caption{Simulation-based recovery across generating angles. Low angles show strong upward bias and poor coverage, whereas moderate and high angles are recovered substantially better.}
\label{fig:supp-recovery-plot}
\end{figure}

\subsection*{Matched-domain calibration in the empirical low-angle regime}

Calibration was evaluated directly at the empirical low-angle solutions observed in the matched Gender--Science and Sexuality datasets. In both domains, recovery remained strongly biased upward, and interval coverage was poor at the relevant sample sizes. Under the present binned Gaussian profiling model and empirical online-IAT noise structure, low-angle profile peaks define a stable population-level descriptor with limited low-angle generative recovery.

In the Gender--Science corpus, calibration at $\theta_{\mathrm{true}}=17.25^\circ$ yielded mean recovered MAP values of $39.04^\circ$ at $N=10{,}000$ and $39.62^\circ$ at $N=30{,}000$, with 94\% HDI coverage equal to 0\% at both sample sizes. In the matched Sexuality corpus, calibration at $\theta_{\mathrm{true}}=19.12^\circ$ showed the same pattern: mean recovered MAP $41.99^\circ$ at $N=10{,}000$ and $38.94^\circ$ at $N=30{,}000$, again with 0\% coverage.

\subsection*{Null controls}

Two null controls were used to assess whether the observed posterior concentration could be explained by disrupted temporal order or generic smooth trends alone.

\begin{figure}[htbp]
\centering
\includegraphics[width=0.95\textwidth]{ext_null_distributions.png}
\caption{Null posterior diagnostics versus the observed fit.}
\label{fig:supp-nulls}
\end{figure}

\begin{table}[htbp]
\centering
\caption{Summary of posterior-shape diagnostics under the two null families.}
\label{tab:supp-null-summary}
\small
\begin{tabular}{@{}lcccccc@{}}
\toprule
Null type & Replicates & MAP mean & MAP SD & Mean posterior SD & Mean HDI width & Mean entropy \\
\midrule
Permutation & 500 & 109.650 & 67.790 & 7.716 & 24.120 & 1.317 \\
Quadratic   & 300 & 42.250 & 13.590 & 11.591 & 41.700 & 5.200 \\
\bottomrule
\end{tabular}
\end{table}

\subsection*{Model comparison details}

Predictive performance was evaluated by participant-wise temporal hold-out, with the highest-$x$ approximately 20\% of bins reserved for testing.

\begin{table}[htbp]
\centering
\caption{Hold-out RMSE by model for the main six-bin analysis.}
\label{tab:supp-model-comparison}
\small
\begin{tabular}{@{}lccccc@{}}
\toprule
Model & Best parameter & RMSE$_{\mathrm{test}}$ & Mean log score$_{\mathrm{test}}$ & $\sigma_{\mathrm{train}}$ & Boundary hit \\
\midrule
Interference & 124.20 & 21.919 & -21524.695 & 0.1056 & No \\
Cosine only & 180.00 & 21.899 & -7212.871 & 0.1823 & Yes \\
Sine only & 180.00 & 21.897 & -9632.215 & 0.1578 & Yes \\
Polynomial (\(2^{\mathrm{nd}}\) order) & --- & 21.924 & -21602.217 & 0.1055 & No \\
Exponential & 10.10 & 21.896 & -9905.413 & 0.1556 & No \\
Power law & 0.05 & 21.897 & -9259.726 & 0.1609 & Yes \\
\bottomrule
\end{tabular}
\end{table}

\subsection*{Age IAT supplementary results}

The Age IAT was added as an additional conventional public social-category IAT. Trial-level data were downloaded automatically from the Project Implicit OSF archive and processed using the same critical-block pipeline applied to the Gender--Science and Sexuality domains. Analyses retained blocks 3, 4, 6, and 7 and used a cleaned latency variable restricted to the 300--10{,}000 ms window.

\begin{table}[htbp]
\centering
\caption{Age IAT posterior summary.}
\label{tab:supp-age-posterior}
\small
\begin{tabular}{@{}lrrrrrr@{}}
\toprule
Domain & Mean $\theta$ & SD $\theta$ & MAP $\theta$ & HDI low & HDI high & Posterior entropy \\
\midrule
Age & 18.476 & 0.522 & 18.500 & 17.500 & 19.250 & 2.153 \\
\bottomrule
\end{tabular}
\end{table}

\begin{table}[htbp]
\centering
\caption{Age IAT quality-control summary after automatic OSF download and trial-level preprocessing.}
\label{tab:age-iat-qc}
\small
\begin{tabular}{@{}lr@{}}
\toprule
Metric & Value \\
\midrule
Raw rows & 30{,}168{,}625 \\
Retained critical-block trials after latency cleaning & 17{,}907{,}400 \\
Participants & 156{,}005 \\
Available blocks & 0--7 \\
Retained critical blocks & 3, 4, 6, 7 \\
Pre-cleaning median latency (ms) & 734.00 \\
Pre-cleaning mean latency (ms) & 1019.13 \\
Pre-cleaning SD latency (ms) & 43{,}934.52 \\
Post-cleaning median latency (ms) & 741.00 \\
Post-cleaning mean latency (ms) & 944.89 \\
Post-cleaning SD latency (ms) & 673.43 \\
Post-cleaning IQR latency (ms) & 430.00 \\
Fast-trial proportion ($<300$ ms) & 0.0189 \\
Long-trial proportion ($>10{,}000$ ms) & 0.0012 \\
Error rate & 0.0727 \\
\bottomrule
\end{tabular}
\end{table}

\subsection*{Matched Sexuality supplementary results}

To complement the six-bin matched-domain comparison in the main text, the full profiling pipeline was repeated at 4 and 8 bins for both the year-matched 2019 Gender--Science and Sexuality archives.

\begin{table}[htbp]
\centering
\caption{Matched 2019 public-domain pooled \(\theta\) summaries by binning resolution.}
\label{tab:supp-matched-bins}
\footnotesize
\resizebox{\textwidth}{!}{\begin{tabular}{lrrrrrrrrr}
\toprule
Domain & Bins & Curves & Mean $\theta$ & SD $\theta$ & MAP $\theta$ & HDI low & HDI high & Entropy & RSS at MAP \\
\midrule
Gender-Science & 4 & 138957 & 22.901 & 0.628 & 22.750 & 21.750 & 24.000 & 2.338 & 108860.878 \\
Gender-Science & 6 & 138955 & 17.294 & 0.488 & 17.250 & 16.500 & 18.250 & 2.086 & 330497.343 \\
Gender-Science & 8 & 138955 & 15.565 & 0.427 & 15.500 & 14.750 & 16.250 & 1.953 & 539470.098 \\
Sexuality & 4 & 207367 & 25.658 & 0.710 & 25.500 & 24.500 & 27.000 & 2.462 & 163468.302 \\
Sexuality & 6 & 207366 & 19.119 & 0.532 & 19.000 & 18.250 & 20.000 & 2.172 & 503741.487 \\
Sexuality & 8 & 207365 & 17.049 & 0.456 & 17.000 & 16.250 & 17.750 & 2.018 & 826124.504 \\
\bottomrule
\end{tabular}
}
\end{table}

\begin{table}[htbp]
\centering
\caption{Posterior contrast at six bins: Sexuality minus Gender--Science.}
\label{tab:supp-matched-diff}
\small
\begin{tabular}{lrrrrrr}
\toprule
contrast & delta_mean & delta_sd & delta_map & hdi_low & hdi_high & prob_gt_zero \\
\midrule
Sexuality - Gender-Science & 1.825 & 0.722 & 1.750 & 0.500 & 3.000 & 0.991 \\
\bottomrule
\end{tabular}

\end{table}

\begin{table}[htbp]
\centering
\caption{Leak-free temporal hold-out model comparison for both matched 2019 public domains.}
\label{tab:supp-matched-model}
\footnotesize
\resizebox{\textwidth}{!}{\begin{tabular}{lrrlrrrrrrrr}
\toprule
Domain & Year & Bins & Model & Best param. & Train RSS & Test RMSE & Test log score & Train $\sigma$ & Curves & Test obs. & Boundary hit \\
\midrule
Gender-Science & 2019 & 6 & Interference & 124.3000 & 6213.0978 & 22.0670 & -21818.2249 & 0.1056 & 138955 & 277910 & False \\
Gender-Science & 2019 & 6 & Cosine only & 180.0000 & 18459.1189 & 22.0477 & -7317.7371 & 0.1822 & 138955 & 277910 & True \\
Gender-Science & 2019 & 6 & Sine only & 180.0000 & 13822.3879 & 22.0453 & -9770.9012 & 0.1577 & 138955 & 277910 & True \\
Gender-Science & 2019 & 6 & Polynomial & NaN & 6230.4583 & 22.0725 & -21896.5903 & 0.1055 & 138955 & 277910 & False \\
Gender-Science & 2019 & 6 & Exponential & 10.1000 & 13442.1618 & 22.0443 & -10046.7433 & 0.1555 & 138955 & 277910 & False \\
Gender-Science & 2019 & 6 & Power law & 0.0500 & 14436.1647 & 22.0454 & -9393.3915 & 0.1608 & 138955 & 277910 & True \\
Sexuality & 2019 & 6 & Interference & 119.8000 & 10741.8804 & 25.0148 & -24196.6325 & 0.1137 & 207366 & 414732 & False \\
Sexuality & 2019 & 6 & Cosine only & 180.0000 & 27003.5082 & 24.9999 & -9598.2930 & 0.1804 & 207366 & 414732 & True \\
Sexuality & 2019 & 6 & Sine only & 180.0000 & 22236.5202 & 24.9970 & -11653.6571 & 0.1637 & 207366 & 414732 & True \\
Sexuality & 2019 & 6 & Polynomial & NaN & 10757.3621 & 25.0192 & -24275.4904 & 0.1135 & 207366 & 414732 & False \\
Sexuality & 2019 & 6 & Exponential & 8.1500 & 21885.1888 & 24.9966 & -11840.5722 & 0.1624 & 207366 & 414732 & False \\
Sexuality & 2019 & 6 & Power law & 0.0700 & 22587.7528 & 24.9977 & -11495.3729 & 0.1649 & 207366 & 414732 & False \\
\bottomrule
\end{tabular}
}
\end{table}

\subsection*{Independent life-satisfaction IAT angular summary}

The independent life-satisfaction IAT was analyzed using the same pooled angular-profiling framework, with the within-block coordinate defined by normalized display order within the critical combined blocks. The posterior was broader than in the large public Project Implicit corpora, consistent with the smaller sample size, but the modal estimate remained in the low-angle range.

\begin{table}[htbp]
\centering
\caption{Independent life-satisfaction IAT posterior summary.}
\label{tab:supp-life-theta}
\small
\begin{tabular}{@{}lrrrrrr@{}}
\toprule
Domain & Participants & Mean $\theta$ & SD $\theta$ & MAP $\theta$ & HDI low & HDI high \\
\midrule
Life-satisfaction IAT & 180 & 28.49 & 13.68 & 17.75 & 8.50 & 53.75 \\
\bottomrule
\end{tabular}
\end{table}

\subsection*{Life-satisfaction IAT sample}

\begin{table}[htbp]
\centering
\caption{Sociodemographic characteristics of the independent life-satisfaction IAT sample.}
\label{tab:supp-life-iat-demographics}
\footnotesize
\begin{tabular}{@{}lcccc@{}}
\toprule
& \multicolumn{2}{c}{Individualist ($n=68$)} & \multicolumn{2}{c}{Collectivist ($n=112$)} \\
\cmidrule(lr){2-3}\cmidrule(lr){4-5}
Characteristic & $n$ & \% & $n$ & \% \\
\midrule
\multicolumn{5}{@{}l}{\textit{Sex at birth}} \\
Female & 43 & 63.23 & 86 & 76.78 \\
Male & 25 & 36.77 & 26 & 23.22 \\[2pt]
\multicolumn{5}{@{}l}{\textit{Marital status}} \\
Single & 45 & 66.18 & 70 & 62.50 \\
Married & 18 & 26.47 & 34 & 30.36 \\
Separated & 4 & 5.88 & 6 & 5.36 \\
Widowed & 1 & 1.47 & 2 & 1.79 \\[2pt]
\multicolumn{5}{@{}l}{\textit{Schooling}} \\
Middle school & 38 & 55.88 & 58 & 51.79 \\
Complete higher education & 15 & 22.06 & 15 & 13.39 \\
Postgraduate studies & 15 & 22.06 & 39 & 34.82 \\[2pt]
\multicolumn{5}{@{}l}{\textit{Ethnicity}} \\
White & 40 & 58.82 & 61 & 54.46 \\
Brown & 23 & 33.82 & 37 & 33.04 \\
Black & 4 & 5.88 & 10 & 8.93 \\
Indigenous & 0 & 0.00 & 2 & 1.79 \\[2pt]
\multicolumn{5}{@{}l}{\textit{Country region}} \\
South & 17 & 25.00 & 39 & 34.82 \\
Southeast & 27 & 39.71 & 32 & 28.57 \\
Midwest & 3 & 4.41 & 3 & 2.68 \\
Northeast & 19 & 27.94 & 33 & 29.46 \\
North & 5 & 7.35 & 5 & 4.46 \\
\bottomrule
\end{tabular}
\end{table}

\end{document}